\documentclass[../../main.tex]{subfiles}
\begin{document}
% SOURCE p=238,239; printed=223-224; object=ch07-vague-distribution; number=7.19; region=bottom-top
\begin{proposition}\label{proposition:ch07-vague-distribution}
设 $\mu,\mu_1,\mu_2,\ldots\in M(\mathbb R)$，令 $F_n(x)=\mu_n(({-\infty},x])$、$F(x)=\mu(({-\infty},x])$。
\begin{enumerate}
\item\label{item:ch07-distribution-implies-vague} 若 $\sup_n\|\mu_n\|<\infty$，且在 $F$ 的每个连续点 $x$ 都有 $F_n(x)\to F(x)$，则 $\mu_n\to\mu$ 模糊收敛。
\item\label{item:ch07-vague-implies-distribution} 若 $\mu_n\geqslant0$，$\lim_{x\to-\infty}[\sup_nF_n(x)]=0$，且 $\mu_n\to\mu$ 模糊收敛，则在 $F$ 的每个连续点 $x$ 都有 $F_n(x)\to F(x)$。
\end{enumerate}
\end{proposition}
% SOURCE p=239; printed=224; object=ch07-vague-distribution-proof; region=upper-middle
\begin{proof}
第一项：由\cref{theorem:ch03-bv-properties,theorem:ch03-nbv-measure-correspondence}，$F$ 除了至多可数个点外处处连续，所以 $F_n\to F$ 对 Lebesgue 测度几乎处处成立。此外，$\|F_n\|_u\leqslant\|\mu_n\|$，故 $F_n$ 一致有界。若 $f$ 连续可微且具有紧支集，由分部积分（\cref{theorem:ch03-stieltjes-parts}）和控制收敛定理，
\begin{align*}
\int f\,\mathrm d\mu_n
=-\int f'(x)F_n(x)\,\mathrm dx
\longrightarrow-\int f'(x)F(x)\,\mathrm dx
=\int f\,\mathrm d\mu.
\end{align*}
这里修正了原书中两个积分前漏印的负号，具体说明见下。由\cref{theorem:ch04-noncompact-stone-weierstrass}，所有这样的 $f$ 在 $C_0(\mathbb R)$ 中稠密；再由\cref{proposition:ch05-strong-dense-convergence}，积分收敛对每个 $f\in C_0(\mathbb R)$ 都成立。因此 $\mu_n\to\mu$ 模糊收敛。

第二项：给定 $a\in\mathbb R$ 及 $\delta,\varepsilon>0$，取足够大的 $N$ 使 $\sup_nF_n(-N)<\delta$。令 $f\in C_c(\mathbb R)$ 在 $[-N,a]$ 上等于 $1$，在 $({-\infty},-N-\varepsilon]$ 及 $[a+\varepsilon,\infty)$ 上等于 $0$，在中间两段上线性。若 $n$ 足够大，使 $\int f\,\mathrm d\mu_n<\int f\,\mathrm d\mu+\delta$，则
\begin{align*}
F_n(a)-\delta&<F_n(a)-F_n(-N)=\mu_n((-N,a])\leqslant\int f\,\mathrm d\mu_n\\
&<\int f\,\mathrm d\mu+\delta\leqslant F(a+\varepsilon)+\delta.
\end{align*}
由于 $\delta$ 任意，$\limsup_{n\to\infty}F_n(a)\leqslant F(a+\varepsilon)$。类似地，考虑在 $[-N+\varepsilon,a-\varepsilon]$ 上为 $1$，在 $({-\infty},-N]$ 与 $[a,\infty)$ 上为 $0$，中间线性的函数，可得 $\liminf_{n\to\infty}F_n(a)\geqslant F(a-\varepsilon)$。这里原书把后一个函数的左侧零区间印作 $({-\infty},N]$，应为 $({-\infty},-N]$。由于 $\varepsilon$ 任意，若 $F$ 在 $a$ 连续，就有 $F_n(a)\to F(a)$。
\end{proof}
% SUPPLEMENT: id=v14-ch07-vague-distribution-details; parent=ch07-vague-distribution; kind=proof-detail+correction+motivation+hypotheses; origin=independent-derivation
\begin{remark}
\textbf{逐步补证及原式订正（重构新增）。}第一项的分部积分在包含 $\operatorname{supp}f$ 的有限区间上使用，边界项因 $f$ 在两端为零而消失，所以必有负号。原书的链式公式是
\begin{align*}
\int f\,\mathrm d\mu_n=\int f'F_n\longrightarrow\int f'F=\int f\,\mathrm d\mu,
\end{align*}
其中两个无负号的等式都不正确。例如 $\mu_n=\mu=\delta_0$，取 $f\in C_c^1(\mathbb R)$ 且 $f(0)=1$，则 $\int f\,\mathrm d\delta_0=1$，而 $\int f'\chi_{[0,\infty)}=-1$。上面的订正使分部积分正确，结论不受影响。若 $M=\sup_n\|\mu_n\|$，则 $|f'F_n|\leqslant M|f'|$；右侧有紧支集且可积，这完整核验了控制收敛的支配条件。

关于稠密性：$C_c^1(\mathbb R)$ 是对共轭封闭的复代数，并在每点不全为零且分离不同点。具体可平移、缩放函数 $b(t)=(1-t^2)^2$（$|t|\leqslant1$，其余为零）得到在指定点为 $1$、在另一个点为零的 $C_c^1$ 函数。非紧 Stone--Weierstrass 定理便适用。对任意 $f\in C_0$、$g\in C_c^1$，
\begin{align*}
|I_{\mu_n}(f)-I_\mu(f)|\leqslant(M+\|\mu\|)\|f-g\|_u+|I_{\mu_n}(g)-I_\mu(g)|,
\end{align*}
这也直接证明了从稠密集推广的步骤。

第二项首先隐含了 $\mu\geqslant0$，应予证明。对每个非负 $g\in C_c$，$I_\mu(g)=\lim_nI_{\mu_n}(g)\geqslant0$，所以限制到 $C_c$ 的泛函为正。正 Riesz 定理给出正 Radon 测度 $\rho$；由于此泛函已有 $\|I_\mu\|$ 的一致范数界，开集质量公式给 $\rho(X)<\infty$。两个积分泛函由稠密性在 $C_0$ 上相同，复 Riesz 的唯一性给 $\mu=\rho\geqslant0$，所以原证中涉及 $F$ 的大小比较合法。

上界构造时还取 $-N<a$，然后取极限给 $\limsup F_n(a)\leqslant F(a+\varepsilon)+2\delta$，再令 $\delta\downarrow0$。为补全下界，固定 $\varepsilon,\delta>0$，增大 $N$ 使 $-N+\varepsilon<a-\varepsilon$ 且 $F(-N+\varepsilon)<\delta$；后一条件由 $\mu$ 有限及左尾集合递减为空得到。取原证订正后的下包函数 $h$。逐点有 $0\leqslant h\leqslant\chi_{({-\infty},a]}$，故
\begin{align*}
\liminf_nF_n(a)&\geqslant\lim_n\int h\,\mathrm d\mu_n=\int h\,\mathrm d\mu\\
&\geqslant\mu((-N+\varepsilon,a-\varepsilon])
=F(a-\varepsilon)-F(-N+\varepsilon)>F(a-\varepsilon)-\delta.
\end{align*}
令 $\delta\downarrow0$ 得下界。原印的 $({-\infty},N]$ 会与函数等于 $1$ 的区间重叠，因而必须订正负号。最后在 $F$ 的连续点令 $\varepsilon\downarrow0$，上下界同趋 $F(a)$。

动机是用连续帽形函数逼近半直线示性函数。第一项的统一总变差界控制稠密逼近误差；连续点条件排除了边界原子的跳跃。第二项的正性使函数的逐点夹逼能传递到积分，统一左尾条件阻止质量逃向 $-\infty$。这两个假设各有独立作用，\cref{exercise:ch07-distribution-assumptions}会给出反例。该命题针对实直线的半直线分布函数，不是任意 LCH 空间上无须修改即可使用的表述。
\end{remark}
% END SUPPLEMENT: id=v14-ch07-vague-distribution-details
\subfile{section-03-exercises.tex}
\end{document}
